\chapter{حقوق زنان و اقلیت‌ها: پایه‌های برابری}
\label{app:rights}

دموکراسی در ایران بدون تأمین حقوق کامل نیمی از جامعه و تکثر قومی و مذهبی کشور صوری خواهد بود.

\section{منشور حقوق زنان در ایران نوین}

طرح مهستان بر لغو تمامی قوانین تبعیض‌آمیز جنسیتی در روز اول گذار تأکید دارد.

\begin{modernbox}{اقدامات فوری برای حقوق زنان}
\begin{enumerate}
    \item \textbf{آزادی پوشش:} لغو هرگونه قانون حجاب اجباری و گشت‌های کنترل اجتماعی.
    \item \textbf{برابری حقوقی:} اصلاح قوانین خانواده، ارث، دیه و حق طلاق بر اساس برابری کامل.
    \item \textbf{مشارکت سیاسی:} اعمال سهمیه حداقل ۳۰٪ در تمام نهادهای دوران گذار.
    \item \textbf{الحاق به CEDAW:} پیوستن به کنوانسیون رفع کلیه اشکال تبعیض علیه زنان بدون قید و شرط.
\end{enumerate}
\end{modernbox}

\section{حقوش اقلیت‌ها و تنوع ملی}

رویکرد مهستان نسبت به تنوع قومی ایران، «وحدت در کثرت» از طریق تمرکززدایی است.

\begin{center}
\begin{tikzpicture}[node distance=2.5cm, auto]
    \tikzstyle{block} = [rectangle, draw, fill=SoftGreen!20, text width=6cm, text centered, rounded corners, minimum height=3em]

    \node [block] (lang) {\textbf{حقوق زبانی} \\ آموزش به زبان مادری و استفاده در نهادهای محلی};
    \node [block, below of=lang] (econ) {\textbf{عدالت توزیعی} \\ رفع تبعیض اقتصادی از مناطق محروم و مرزی};
    \node [block, below of=econ] (pol) {\textbf{مشارکت سیاسی} \\ کرسی‌های تضمینی در مجلس و مدیریت محلی};

    \path [draw, line width=1pt] (lang) -- (econ);
    \path [draw, line width=1pt] (econ) -- (pol);
\end{tikzpicture}
\end{center}

\section{جدول مقایسه‌ای رویکرد به اقلیت‌ها}

\begin{center}
\begin{longtable}{R{4cm} R{5cm} R{6cm}}
\caption{تحول در حقوق اقلیت‌ها} \\
\hline
\textbf{شاخص} & \textbf{رویکرد فعلی} & \textbf{رویکرد مهستان} \\
\hline
\endfirsthead
\hline
\textbf{شاخص} & \textbf{رویکرد فعلی} & \textbf{رویکرد مهستان} \\
\hline
\endhead
\hline
\foot
تمرکز قدرت & مرکزگرایی شدید & فدرالیسم / تمرکززدایی گسترده \\
زبان مادری & محدودیت و جرم‌انگاری & آموزش رسمی و ترویج فرهنگی \\
توسعه اقتصادی & تبعیض آمیز (حاشیه vs مرکز) & بودجه‌ریزی عادلانه و جبرانی \\
نمایندگی & گزینشی و نظارتی & انتخابی و تضمینی (۲۰ کرسی اختصاصی) \\
\hline
\end{longtable}
\end{center}
